\documentclass[]{article}

\usepackage{amsmath}
\usepackage{multirow}
\usepackage{natbib}
\usepackage{graphicx} 


\begin{document}

In this paper, we investigated the connection between deterministic Hamiltonian dynamics and emergent anomalous transport statistics. We addressed the question of when and why the complex transport of passive tracers in a chaotic flow deviates from the predictions of a simple, memoryless stochastic model. To this end, we performed numerical simulations of tracer advection in a two-dimensional point-vortex system, a canonical model for Hamiltonian chaos, and systematically compared its transport properties against a L\'evy walk null model.

Our methodology involved simulating tracer trajectories across a range of vortex densities ($N \in \{5, 10, 20, 40\}$) and analyzing the resulting statistics, including the mean squared displacement (MSD), velocity autocorrelation function (VACF), and residence time distributions. A key aspect of our analysis was the rigorous treatment of measurement noise, which allowed us to distinguish genuine physical dynamics from statistical artifacts, particularly in sparse vortex configurations.

Our results reveal a non-monotonic relationship between vortex density and the nature of the transport. For sparse to intermediate densities ($N \le 20$), the system exhibits near-ballistic motion, with an anomalous diffusion exponent $\alpha \approx 2$. In this regime, tracers are advected coherently by the large-scale flow structures. However, a distinct transition occurs in the densest configuration ($N=40$), where the transport becomes markedly superdiffusive, with $\alpha \approx 1.68$. This transition is driven by a change in the underlying dynamics. Mechanistic analysis demonstrated that the point-vortex system possesses significant Hamiltonian memory, evidenced by persistent velocity correlations that are, by construction, absent in the memoryless L\'evy walk model. Furthermore, we found that increasing vortex density leads to heavier-tailed residence time distributions, indicating that tracers in more chaotic environments are subject to more frequent and prolonged trapping events.

From these findings, we have learned that vortex density acts as a critical control parameter that governs a transition between distinct anomalous transport regimes. The shift from near-ballistic to superdiffusive motion at high density is not a simple loss of transport efficiency but is mechanistically linked to the interplay between trapping and flight dynamics. The increased chaotic interactions in the dense system disrupt long, coherent flights and enhance tracer trapping, a memory-dependent process. This work establishes that while a simple stochastic model like a L\'evy walk may replicate certain macroscopic features, such as the scaling of the MSD, it fails to capture the essential role of deterministic memory in shaping the transport dynamics of Hamiltonian systems. Our findings underscore the importance of considering these memory effects for a complete physical description of transport in chaotic flows.

\end{document}
                